\documentclass[12pt,a4paper]{article}
\usepackage[utf8]{inputenc}
\usepackage[T1]{fontenc}
\usepackage{amsmath,amssymb,amsthm,mathrsfs,mathtools}
\usepackage[margin=1in]{geometry}
\usepackage{hyperref}
\hypersetup{colorlinks=true,linkcolor=blue,citecolor=blue,urlcolor=blue}
\usepackage{booktabs}

\theoremstyle{plain}
\newtheorem{theorem}{Theorem}[section]
\newtheorem{proposition}[theorem]{Proposition}
\newtheorem{lemma}[theorem]{Lemma}
\newtheorem{corollary}[theorem]{Corollary}
\theoremstyle{definition}
\newtheorem{definition}[theorem]{Definition}
\theoremstyle{remark}
\newtheorem{remark}[theorem]{Remark}

\newcommand{\R}{\mathbb{R}}
\newcommand{\C}{\mathbb{C}}
\newcommand{\dd}{\mathrm{d}}
\newcommand{\cM}{\mathcal{M}}

\title{The Spectral Weight $\pi\lambda/\sinh\pi\lambda$:\\[4pt]
Derivation from Phase Space, Modular Form, and Exact Moments}

\author{Daniel Toupin\\[4pt]
\textit{Golden Physics Project}\\[2pt]
\texttt{goldenphysics.org}}

\date{August 2026 --- v2, superseding \texttt{haar\_qg\_paper\_v2.1.5.1}
}

\begin{document}
\maketitle

\begin{abstract}
\noindent
The function $P(\lambda)=\pi\lambda/\sinh(\pi\lambda)$ has appeared
throughout this series as a spectral weight of unclear provenance. Its
origin is now settled. In a companion paper the two-particle massless
unitarity cut is computed in celestial variables, and its Mellin image
is $\Gamma(\Delta_5)\Gamma(\Delta_6)/\Gamma(\Delta_5+\Delta_6)$ times
elementary factors; restricted to the shadow locus
$\Delta_5+\Delta_6=2$ with $\Delta_5=1+i\lambda$, this equals
$\Gamma(1+i\lambda)\Gamma(1-i\lambda)=P(\lambda)$. So $P$ is a
phase-space factor of the cut. It is a derived quantity, and we prove
here the identities it satisfies: a Planck form
$P=2\pi\lambda[n_B(\pi\lambda)-n_B(2\pi\lambda)]$ with the zero-point
energies cancelling, an elementary Fourier pair with
$\tfrac14\operatorname{sech}^2(x/2)$, and a ladder-moment sequence whose
first two members are $\cM_1=1/8$ and $\cM_2=1/90=\zeta(4)/\pi^4$,
verified independently here to 26 digits.

This paper also retracts, in full and by name, the physical claims
built on this function in the draft it replaces. The parallel
retractions for the blackbody draft are in \cite{ToupinModular}. $P$ is not the
Plancherel density of $\mathrm{SL}(2,\C)$, which grows like $\lambda^2$
for the scalar principal series. $P$ does not replace the Feynman loop
measure. The bound $\cM_L\le(1/8)^L$ is a true statement about the
convolution integrals $\cM_L$ and a false statement about loop
amplitudes, so the claim of ultraviolet finiteness without
counterterms does not stand: the companion derivation reconstructs the
standard loop integral exactly, divergences included. The fifteen-digit
box verification reported in the earlier draft compared a
Feynman-parameter integral with a reparameterization of itself. What
survives is a family of exact special-function identities, and a
physical interpretation of $P$ as the spectral density of a
two-particle cut rather than as a cure for anything.
\end{abstract}

\medskip
\noindent\textbf{Keywords:} celestial holography, unitarity, Plancherel
measure, modular theory, Bisognano--Wichmann, Bernoulli numbers.

\tableofcontents
\newpage

% =====================================================================
\section{Where the weight comes from}
% =====================================================================

Earlier drafts in this series postulated $P(\lambda)$ as a loop measure
and then looked for reasons it should be one. The order of reasoning
was backwards, which is how the errors listed in
Sec.~\ref{sec:retractions} entered. The correct order runs as follows.

\begin{theorem}[Phase-space origin]
\label{thm:origin}
Let $\dd\Pi_2$ be two-particle massless phase space at total mass $M$,
written in celestial coordinates. Then
\begin{equation}
\int \dd\Pi_2\;\omega_5^{\Delta_5-1}\omega_6^{\Delta_6-1}
=\frac{1}{8\pi}\Bigl(\frac{M}{2}\Bigr)^{\Delta_5+\Delta_6-2}
\frac{\Gamma(\Delta_5)\Gamma(\Delta_6)}{\Gamma(\Delta_5+\Delta_6)} ,
\label{eq:phi}
\end{equation}
and on the shadow locus $\Delta_5+\Delta_6=2$ with
$\Delta_5=1+i\lambda$,
\begin{equation}
\frac{\Gamma(\Delta_5)\Gamma(\Delta_6)}{\Gamma(\Delta_5+\Delta_6)}
=\Gamma(1+i\lambda)\Gamma(1-i\lambda)
=\frac{\pi\lambda}{\sinh(\pi\lambda)}
=:P(\lambda).
\label{eq:origin}
\end{equation}
\end{theorem}

\begin{proof}
Equation \eqref{eq:phi} is Theorem 4.1 of \cite{ToupinLoops}, proved
there from the celestial parameterization with all Jacobians computed,
and verified to between $10^{-20}$ and $10^{-31}$ at five points with
complex $\Delta$. Equation \eqref{eq:origin} then follows from
$\Gamma(2)=1$ and the reflection formula, checked here to $10^{-41}$.
\end{proof}

The locus $\Delta_5+\Delta_6=2$ is where the $M$ dependence of
\eqref{eq:phi} drops out, so it is the locus on which the cut is scale
invariant. Nothing is localized there by a pole, and no residue is
taken. This corrects the mechanism assumed throughout the earlier
drafts.

% =====================================================================
\section{Identities satisfied by the weight}
% =====================================================================

\begin{theorem}[Planck form, with cancelling zero-point energy]
\label{thm:planck}
With $n_B(y)=(e^y-1)^{-1}$ and $E(x)=\tfrac{x}{2}\coth\tfrac{x}{2}$,
\begin{equation}
P(\lambda)=2\pi\lambda\bigl[n_B(\pi\lambda)-n_B(2\pi\lambda)\bigr]
=2\pi\lambda\sum_{n\ge0}e^{-2\pi\lambda(n+\frac12)},
\qquad
2E(\pi\lambda)-E(2\pi\lambda)=P(\lambda).
\label{eq:planck}
\end{equation}
The zero-point terms cancel identically:
$2\cdot\tfrac{\pi\lambda}{2}-\tfrac{2\pi\lambda}{2}=0$.
\end{theorem}

\begin{proof}
Both statements reduce to
$\coth y-\coth 2y=1/\sinh 2y$ after clearing factors. Verified at three
values of $\lambda$ to $3\times10^{-39}$ or better.
\end{proof}

\begin{proposition}[Fourier pair]
\label{prop:fourier}
\begin{equation}
P(\lambda)=\int_\R e^{i\lambda x}\frac{\dd x}{4\cosh^2(x/2)},
\qquad
\frac{1}{2\pi}\int_\R P(\lambda)\cos(\lambda y)\dd\lambda
=\frac{1}{4\cosh^2(y/2)} .
\label{eq:fourier}
\end{equation}
Verified at two values to $1.7\times10^{-41}$.
\end{proposition}

\begin{proposition}[Logarithmic expansion]
\label{prop:log}
For $|\lambda|<1$,
$\log P(\lambda)=-\sum_{k\ge1}\frac{(-1)^{k+1}\zeta(2k)}{k}\lambda^{2k}$,
so the Taylor coefficients of $\log P$ are even zeta values.
\end{proposition}

\begin{remark}[Thermal reading, treated separately]
\label{rem:thermal}
Equation \eqref{eq:planck} admits a reading through the
Bisognano--Wichmann theorem, in which $P$ is the partition function of
one oscillator at the modular temperature $\beta=2\pi$ weighted by
boost frequency. That structure, together with the Matsubara pole
positions, the Weierstrass product, and the Mellin moment family
$\mathfrak m_s=\pi^{-(s+1)}(1-2^{-(s+1)})\Gamma(s+1)\zeta(s+1)$, is
developed in \cite{ToupinThermal}. The ladder moment $\cM_1$ below is
the $s=1$ member of that family.
\end{remark}

% =====================================================================
\section{Ladder moments}
% =====================================================================

\begin{definition}[Ladder moments]
\label{def:moments}
\begin{equation}
\cM_L:=\frac{1}{(2\pi)^L}\int_{\R_{>0}^L}
\Bigl[\prod_{j=1}^{L}P(\lambda_j)\Bigr]
\Bigl[\prod_{j=1}^{L-1}P\bigl(|\lambda_j-\lambda_{j+1}|\bigr)\Bigr]
\dd\lambda_1\cdots\dd\lambda_L ,
\label{eq:moments}
\end{equation}
with the empty product equal to $1$ at $L=1$.
\end{definition}

\begin{theorem}[First two moments]
\label{thm:moments}
$\cM_1=\tfrac18$ and $\cM_2=\tfrac{1}{90}=\zeta(4)/\pi^4$.
\end{theorem}

\begin{proof}
For $L=1$, $\int_0^\infty \pi\lambda/\sinh(\pi\lambda)\dd\lambda=\pi/4$
by the classical value $\int_0^\infty x/\sinh x\,\dd x=\pi^2/4$, so
$\cM_1=(\pi/4)/(2\pi)=1/8$, exactly. For $L=2$ the double integral was
recomputed here by independent quadrature, giving
$0.011111111111111111$ against $1/90$ with relative error
$1.8\times10^{-26}$.
\end{proof}

\begin{remark}[Status of these numbers]
Definition \ref{def:moments} is a definition of convolution integrals
of $P$. Theorem \ref{thm:moments} is a fact about those integrals.
Neither statement asserts that $\cM_L$ normalizes an $L$-loop
amplitude. The earlier draft made that assertion, and it is withdrawn.
The bound $\cM_L\le(1/8)^L$, which follows from
$0<P\le1$ applied to the rung factors, remains true of
\eqref{eq:moments} and says nothing about ultraviolet behaviour of
amplitudes.
\end{remark}

% =====================================================================
\section{Retractions}
\label{sec:retractions}
% =====================================================================

The following claims from
\texttt{haar\_qg\_paper\_v2.1.5.1} are withdrawn. They are listed
individually so that anyone holding those drafts can mark them.

\begin{enumerate}
\item \textbf{``$P(\lambda)$ is the Plancherel measure of
$\mathrm{SL}(2,\C)$.''} False. The Plancherel density for the scalar
principal series is proportional to $\lambda^2$. The correct
description of $P$ is Theorem \ref{thm:origin}: a phase-space factor of
the two-particle cut, equal to a ratio of Gamma functions on the shadow
locus.

\item \textbf{``The flat Feynman spectral weight is replaced by
$P(\lambda)$.''} Withdrawn. The companion derivation
\cite{ToupinLoops} reconstructs the ordinary loop integral from its cut
by an unsubtracted dispersion relation, verified to relative error
below $1.6\times10^{-21}$. No measure is replaced anywhere in that
chain.

\item \textbf{``$\cM_L\le(1/8)^L$, so the loop expansion is ultraviolet
finite with no regularization and no counterterms.''} The inequality is
true of Definition \ref{def:moments}. The physical conclusion does not
follow and is withdrawn. Divergences of the box appear exactly where
standard methods put them; in \cite{ToupinLoops} they are handled by an
internal mass regulator, with the collinear logarithm confined to a
subleading Mellin pole.

\item \textbf{``$\cM_1=1/8$ and $\cM_2=1/90$ are the one- and two-loop
measures.''} They are values of the integrals \eqref{eq:moments}, and
the rational values are a coincidence of the ladder graph topology at
$L\le2$, as the earlier draft itself argued. The identification with
loop normalizations is withdrawn.

\item \textbf{``The one-loop box is verified to fifteen significant
figures with a 19.6-fold speedup.''} Withdrawn. The routine
\texttt{box\_shadow\_1d} in \texttt{verify\_qg\_measure.py} integrates
one Feynman parameter analytically and maps a second by a logistic
change of variable, so \texttt{check\_box} compares a Feynman-parameter
integral against a reparameterization of itself. The comparison cannot
test any celestial construction. That routine is deprecated and should
not be cited.

\item \textbf{``The ultraviolet catastrophe of cavity radiation and the
ultraviolet divergence of quantum gravity are the same disease with the
same cure.''} Withdrawn as physics, and separately withdrawn in
\cite{ToupinThermal}. The Planck-form identity \eqref{eq:planck} and
the thermal reading of Remark \ref{rem:thermal} stand as exact
statements about the function $P$.

\item \textbf{Transcendentality claims tied to the digamma moment.} The
identity
$\tfrac{1}{2\pi}\int_\R P\,\mathrm{Re}\,\psi(\tfrac12+\tfrac{i\lambda}{2})
\dd\lambda=\tfrac18+\tfrac14\psi(\tfrac12)$ is proved in
\cite{ToupinBlocks}. The $\tfrac18$ there arises from a full-line
integral, while $\cM_1=\tfrac18$ is a half-line integral, so reading
the moment as ``the one-loop measure plus a quarter of the digamma''
overstates what has been shown.
\end{enumerate}

% =====================================================================
\section{What stands}
% =====================================================================

\begin{center}
\begin{tabular}{lll}
\toprule
Statement & Status & Check \\
\midrule
$P=\Gamma(1+i\lambda)\Gamma(1-i\lambda)$ from phase space & proved & $10^{-41}$ \\
Planck form, zero-point cancellation & proved & $3\times10^{-39}$ \\
Fourier pair with $\tfrac14\operatorname{sech}^2(x/2)$ & proved & $1.7\times10^{-41}$ \\
$\log P$ coefficients are even zeta values & proved & classical \\
Planck form and thermal reading & see \cite{ToupinThermal} & \\
$\cM_1=1/8$ & proved & exact \\
$\cM_2=1/90=\zeta(4)/\pi^4$ & verified & $1.8\times10^{-26}$ \\
$P$ as Plancherel density & \textbf{withdrawn} & false \\
$P$ as loop-measure replacement & \textbf{withdrawn} & false \\
UV finiteness without counterterms & \textbf{withdrawn} & false \\
Fifteen-digit box check & \textbf{withdrawn} & circular \\
\bottomrule
\end{tabular}
\end{center}

\begin{thebibliography}{9}
\bibitem{ToupinLoops} D.~Toupin, ``Loop amplitudes from unitarity cuts
in celestial Mellin space: a complete derivation for the scalar box,''
v19, 2026.
\bibitem{ToupinBlocks} D.~Toupin, ``Principal-series blocks on the
celestial sphere,'' v2, 2026.
\bibitem{BW1975} J.~Bisognano, E.~Wichmann, J.\ Math.\ Phys.\ 16 (1975)
985.
\bibitem{BW1976} J.~Bisognano, E.~Wichmann, J.\ Math.\ Phys.\ 17 (1976)
303.
\bibitem{ToupinThermal} D.~Toupin, ``Modular thermality of the
celestial spectral weight,'' v2, 2026.
\bibitem{ToupinModular} D.~Toupin, ``Modular thermality of the
celestial spectral weight,'' v2, 2026.
\bibitem{SimmonsDuffin2014} D.~Simmons-Duffin, ``TASI lectures on the
conformal bootstrap,'' arXiv:1602.07982.
\end{thebibliography}

\end{document}
